% celeste-sat technical report -- DRAFT.
% Build: cd docs/report && pdflatex report && bibtex report && pdflatex report && pdflatex report
% \TODO{} marks what is still missing (author details, citations). The Chapter 1 numbers are from results/certificates/lvl_1_w8.json.
\documentclass[conference]{IEEEtran}
\usepackage[T1]{fontenc}
\usepackage{booktabs}
\usepackage{xcolor}
\usepackage{url}
\usepackage{amsmath}
\newcommand{\TODO}[1]{\textcolor{red}{[TODO: #1]}}
\newcommand{\code}[1]{\texttt{#1}}
% filled in from results/certificates/ledge_w8.json
\newcommand{\LEDGEROWS}{all 9 windows ($K = 0$ to 8) are \emph{no} for CBMC, UNSAT for CaDiCaL and verified by drat-trim}
\newcommand{\LEDGECHECK}{4 to 19 minutes}

\begin{document}

\title{Proving a Celeste Speedrun Locally Optimal:\\
A Bit-Exact Model, SAT Queries and Checked Certificates}

\author{\IEEEauthorblockN{\TODO{author}}
\IEEEauthorblockA{\TODO{affiliation, e-mail}}}

\maketitle

\begin{abstract}
We ask whether the community tool-assisted speedrun (TAS) of Chapter~1 of the
platformer \emph{Celeste} can be made faster. We port the game's movement
code to C, check it against the real game frame by frame down to the
subpixel, and use the same code three ways: as a simulator, as a beam
search, and as a CBMC harness in which every frame's input is free. SAT
queries then either find a faster stretch of route or prove that none exists
within the model. On a supercomputer campaign of about 2,300 core-hours,
3,899 of 3,920 windows of 12, 16 and 18 frames over the chapter were proven
impossible to shorten to the same state, and all 36 queries around the 18 room
transitions were proven. To remove the SAT solver from the trusted base, we
re-solve windows with CaDiCaL, emit DRAT proofs, and check them with
drat-trim. What remains trusted is the translation of the harness to CNF and the model itself.
\end{abstract}

\section{The Question}
Speedrunners and TAS authors optimise routes through \emph{Celeste} by hand,
frame by frame. The community TAS of Chapter~1~\cite{celestetas} takes 2,887
frames (0:49.079), of which 2,275 are frames of control in 19 rooms. Is it
optimal?

A game is a good testbed for bounded model checking because its physics is
discrete and bit-precise: positions, speeds and timers are single-precision
floats updated once per frame of $1/60$\,s, and there is ground truth to
check any claim against, namely the game itself. The question is also hard. A room
takes about 100 frames, each frame has about 70 distinct useful inputs, and
the cost of a proof grows roughly $1.5\times$ per free frame. We therefore
answer local questions: can any stretch of $W{+}1$ frames of the route be
done in $W$ frames and reach the same state?

\section{The Model and How It Is Checked}
The model (\code{model/celeste.c}) is a port of the game's
\code{Player.cs} and of the parts of the Monocle engine that move actors:
running, jumps with coyote time and variable height, the 8-direction dash
with its freeze frames, supers, hypers, wall jumps and wallbounces, climbing
and stamina, the 5-frame input buffers, corner correction, spikes,
jump-through platforms, springs, refills, zip movers, falling, crumble and
dash blocks, and room transitions. It reproduces the game's single-precision
arithmetic, including the order of operations. A \emph{literal} build keeps
the game's float timers and pixel-by-pixel movement; an \emph{optimised}
build, which the SAT harness uses, replaces them with integer frame counters and
closed forms.

Fidelity is checked three ways. (i)~Sixteen recordings of Chapter~1, made
in Celeste~1.4 with CelesteTAS, match the model on every frame, subpixels
included; the longest is 880 frames over 8 rooms and ends in a death on the
same frame. (ii)~The whole community TAS replays in the model: every room's
inputs leave the room on exactly their last frame, 2,275 frames in 19 rooms.
(iii)~The literal and the optimised build agree bit for bit on 2.95 million
fuzzed frames in random rooms, and so does every shortcut the harness takes.
Chapter~2, which adds dream blocks, chasers and switch gates, served as a
test: its first recordings exposed four model bugs, all since fixed.

\section{The Queries}
The same C code compiles into a CBMC~\cite{cbmc} harness whose inputs are
nondeterministic. CBMC unrolls $W$ frames into a bit-precise formula and
asks a SAT solver for a counterexample to an assertion.

\emph{Room solve.} ``She does not leave the room within $N$ frames'':
UNSAT proves $N$ too short.
\emph{Windows} (\code{harness/window.c}). The run starts from the TAS's exact
state after $K$ frames; with $W$ free frames, ``she never reaches the TAS's
state after $K{+}W{+}1$ frames, nor the exit''. A counterexample is a route one frame faster; UNSAT
proves that no input sequence gets from the state at $K$ to the state at
$K{+}W{+}1$ in $W$ frames or fewer.
\emph{Crossings} (\code{harness/cross.c}). The same question across a room
transition: the last frames of one room and the first of the next are free.

States are compared on the fields that can still change the rest of the
route. Which fields are dead (e.g.\ a dash cooldown that runs out before the
next dash) is found by replaying the rest of the route with each field
changed (\code{sim/live.c}): 9--14 of 17 candidate fields in room~1. This is
a test, not a proof; a route found with fields left out is replayed, and the
window is asked again comparing every field if the replay is not faster.
Symmetry-breaking rules on inputs (presses that only fill a buffer) are
switched off in a window's last five frames.

\section{Results}
\emph{Faster stretches.} SAT found a frame the human TAS left: room~3
(\code{lvl\_3}) can be left on frame 106 instead of 107 with a wallbounce on
the dash's last frame. The game confirms it, but the next room's inputs fail
from that exit state.

\emph{Window proofs.} On CINECA's Leonardo (about 2,300 core-hours,
CBMC~5.80.0), Table~\ref{tab:leonardo}: 3,899 windows were proven. All 15
``faster'' answers are room \emph{endings}: the window counts leaving the
room as success, and each leaves 2--5 frames sooner in a different state,
after which the next room's TAS inputs fail. The 6 timeouts are not
counted as proven. All 36 crossing queries (the last 6 frames of a room and
3 of the next free, and again 10 and 8) were proven. The campaign ran before the model
had the tiles the game copies above a room's top edge, so the 86 frames near
room edges need re-checking.

\begin{table}[t]
\centering
\caption{Window proofs over the community TAS (Leonardo)}
\label{tab:leonardo}
\begin{tabular}{@{}lrrrr@{}}
\toprule
$W$, rooms & windows & proven & timeout & ``faster'' \\
\midrule
12, all 19 & 2,028 & 2,022 & 2 & 4 \\
16, all 19 & 1,819 & 1,804 & 4 & 11 \\
18, \code{lvl\_1} & 73 & 73 & 0 & 0 \\
\midrule
total & 3,920 & 3,899 & 6 & 15 \\
\bottomrule
\end{tabular}
\end{table}

\emph{Checked certificates.} A ``no'' from CBMC means trusting CBMC's
solver run. Following Seidl's ``never trust your solver''~\cite{seidl}, \code{tools/certify.py}
asks CBMC for the window's formula in DIMACS instead (\code{--dimacs}), solves
it with CaDiCaL~\cite{cadical}, which writes a DRAT proof~\cite{drat}, and
checks the proof with drat-trim~\cite{drattrim}. CBMC's own verdict and
CaDiCaL's must agree, or the run stops. The pipeline is first run where the
answers are known: a demo room with a wasted frame, whose first window must
come back SAT; the next window, which must certify; and that window's proof
with its last tenth cut off, which drat-trim must reject. It does all three.

On the demo room \code{ledge} (an 18-frame route), every window of
$W{=}8$ certifies: \LEDGEROWS. A window's CNF has 2.1 to 3.6 million
variables and 6.5 to 11.1 million clauses; CaDiCaL takes 3 to 9 minutes and
writes a binary DRAT proof of 0.56 to 1.2\,GB, which drat-trim checks in
\LEDGECHECK, about twice the solving time (CBMC~5.95.1, CaDiCaL~1.7.3, one
4-core machine, 8.5\,GB of proofs in all). The CNF that CBMC writes is
deterministic: a window solved by hand and again by the script gives the
same SHA-256.
For Chapter~1's room~1 (\code{lvl\_1}, 92 frames) at $W{=}8$, 81 of the 83
windows certify. Their formulas have 0.3 to 3.9 million variables and up to
11.8 million clauses, the largest proof is 1.5\,GB (76\,GB over all windows,
deleted once checked), and CBMC, CaDiCaL and drat-trim take 16.6~hours in
total. The remaining two windows ($K = 31, 32$) lie entirely inside a dash,
so the inputs have no effect. CBMC proves them during simplification and
emits an empty formula; with nothing left for the SAT solver, there is no
proof, and these two results depend on CBMC alone.

\section{What ``Proven'' Means, and What Is Open}
Every claim is about the model and is usually conditional on the frames
before the window. With certificates, the solver is no longer trusted. What
stays trusted is:
(i)~CBMC's translation of the harness into CNF (unrolling, bit-precise
float arithmetic, the assumptions on inputs), now the largest piece;
(ii)~the model, which is validated by recordings and fuzzing, not proven equal to the game;
(iii)~the question itself (``the same state, one frame sooner'');
(iv)~the field mask, which comes from a test;
(v)~drat-trim, unless replaced by a formally verified checker of LRAT
proofs such as cake\_lpr~\cite{cakelpr}.
Certificates move the trust from the solver to the encoding and the model;
the encoding is the next thing to check.

Two problems are open. The \emph{seams} between rooms: faster endings that
leave a room in a different state might pay off if the next room is
re-routed, which local windows cannot see. And \emph{dominance between
states}: a proof that a different state is no better than the TAS's
(e.g.\ the same position with other subpixels) would turn ``no faster route to
the same state'' into ``no faster route''. Room optimality needs lower bounds,
and run optimality is out of reach for direct unrolling.

\section{Related Work}
Bounded model checking with CBMC~\cite{cbmc}; clausal proofs in DRAT and
their checking with drat-trim~\cite{drat,drattrim}; verified checking of
LRAT proofs~\cite{cakelpr}; the argument for checking a solver's answers
rather than trusting it~\cite{seidl}. The community TAS and its tooling come
from the CelesteTAS project~\cite{celestetas}.
\TODO{Read every cited paper before submission and add only the ones read;
check each entry's details.}

\section*{Acknowledgement}
The code was developed with AI coding agents. Computing time on Leonardo was
provided by EuroHPC (project EUHPC\_D30\_031).

\bibliographystyle{IEEEtran}
\bibliography{report}

\end{document}
